% TOP_PROBLEM: {"id":30007652,"problem_number":"LOCAL-30007652","title":"Identifying genuine poverty traps","category_id":21,"category":{"id":21,"name":"economics","display_name":"Economics","description":"Open economic theory statements, published research agendas, and proposed scoped specifications; question provenance and historical priority are qualified per record.","slug":"economics","order_index":21,"created_at":"2026-10-08T00:00:00Z"},"status":"research_question","external_url":null,"legacy_ids":[],"published":false,"statement":"In the explicitly specified setting: Where do self-reinforcing poverty traps with meaningful asset thresholds exist, and how can they be distinguished from smooth persistence, temporary shocks, and measurement error? The mathematical statement above fixes the target, assumptions and scope of this catalogue formulation.","economics":{"original_id":141,"field":"Development and poverty","kind":"identification","formalization_basis":"adapted_research_question","source_status":"research_agenda","priority_status":"earliest_found","priority_note":"Earliest directly inspected qualifying passage in this audit; earlier origins were not established. A review or research-agenda statement does not imply original priority.","earliest_verified_reference":{"authors":["Dean Karlan","Amol Singh Raswan","Christopher Udry"],"title":"The Sisyphean Pursuit of Evidence for Poverty Traps","year":2026,"url":"https://www.ipr.northwestern.edu/our-work/working-papers/2026/wp-26-24.html","doi":"10.3386/w35253","locator":"Abstract, final paragraph","evidence_summary":"The authors question whether a common population asset threshold can be evidenced under heterogeneity; this is a methodological challenge, not a trap-existence theorem."},"current_reference":{"authors":["Dean Karlan","Amol Singh Raswan","Christopher Udry"],"title":"The Sisyphean Pursuit of Evidence for Poverty Traps","year":2026,"url":"https://www.ipr.northwestern.edu/our-work/working-papers/2026/wp-26-24.html","doi":"10.3386/w35253","locator":"Abstract, final paragraph","evidence_summary":"The authors question whether a common population asset threshold can be evidenced under heterogeneity; this is a methodological challenge, not a trap-existence theorem."},"formalization_note":"The cited passage establishes a research direction, not this exact mathematical statement. The notation, population, policy menu, assumptions, and estimands are an original adaptation for this catalogue; no universal unsolved theorem or source priority is claimed."},"statusQualification":"Published question or research agenda verified at the cited locator. The mathematical specification is adapted for this catalogue and includes additional assumptions; source publication does not establish first-ever priority or certify this exact model remains unsolved. The cited passage establishes a research direction, not this exact mathematical statement. The notation, population, policy menu, assumptions, and estimands are an original adaptation for this catalogue; no universal unsolved theorem or source priority is claimed.","statusReviewedAt":"2026-10-08"}
% FIRST_POSTED: 2026-10-08
% ENTRY_KIND: statement_only
% !TeX program = lualatex
\documentclass[11pt]{article}
\usepackage{fontspec}
\usepackage[a4paper,margin=25mm]{geometry}
\usepackage{amsmath,amssymb,mathtools}
\usepackage{xurl}
\usepackage[hidelinks]{hyperref}
\newcommand{\sref}[2]{\hyperlink{source:#1:#2}{[S#2]}}
\setlength{\emergencystretch}{3em}
\begin{document}
% problemId:30007652
\section*{Identifying genuine poverty traps}\label{30007652}
\subsection{Definitions and mathematical statement}
Fix a population of households, an institutional environment, and a horizon long enough to distinguish recovery from transitory dynamics. Household \(i\)'s productive assets follow \(a_{i,t+1}=f(a_{it},\theta_i,\eta_{it})\), where \(\theta_i\) is persistent heterogeneity and \(\eta_{it}\) a shock with specified distribution. For a type \(\theta\), define expected drift \(d(a,\theta)=E[f(a,\theta,\eta)-a\mid\theta]\). An asset-threshold trap in this particular model requires at least two attracting long-run regimes and an intervening unstable threshold, rather than simply low assets or slow growth. State the relevant stochastic stability definition when shocks permit crossings. The identification task is to distinguish such type-specific nonconvex dynamics from a mixture of households with different smooth dynamics, measurement error, and temporary shocks. Use randomized asset increments over adequate support, repeated measurements, and exogenous shocks to identify or bound \(f\) and long-run responses. Determine whether any common population threshold is empirically defensible. Do not impose it because a fitted aggregate curve appears nonlinear. The cited methodological critique questions whether broad threshold evidence is attainable under heterogeneity, so the task is scoped identification and model discrimination rather than an asserted universal poverty-trap conjecture.

\par\textbf{Formulation and scope.} The cited passage establishes a research direction, not this exact mathematical statement. The notation, population, policy menu, assumptions, and estimands are an original adaptation for this catalogue; no universal unsolved theorem or source priority is claimed.

\par\textbf{Open-question provenance.} An explicit question or research agenda was verified at the source locator. This does not by itself certify that every later version remains unresolved \sref{30007652}{1}.

\par\textbf{Historical priority.} Earliest directly inspected qualifying passage in this audit; earlier origins were not established. A review or research-agenda statement does not imply original priority.

\subsection{Short English statement}
In the explicitly specified setting: Where do self-reinforcing poverty traps with meaningful asset thresholds exist, and how can they be distinguished from smooth persistence, temporary shocks, and measurement error? The mathematical statement above fixes the target, assumptions and scope of this catalogue formulation.

\subsection{Sources}
\begin{itemize}\raggedright
\item[\textnormal{[S1]}]\hypertarget{source:30007652:1}{} Dean Karlan, Amol Singh Raswan, Christopher Udry. 2026. The Sisyphean Pursuit of Evidence for Poverty Traps. Abstract, final paragraph.
\url{https://www.ipr.northwestern.edu/our-work/working-papers/2026/wp-26-24.html}.
DOI: 10.3386/w35253.
{\small\textit{Earliest verified question/agenda reference; historical priority qualified below. The authors question whether a common population asset threshold can be evidenced under heterogeneity; this is a methodological challenge, not a trap-existence theorem.}}
\end{itemize}
\end{document}
